\documentclass[a4paper, 10pt, ]{article}

\input{./COMMONFILES/preamble.tex}
\usepackage{siunitx}

% -----------------------------------------------------------------------------

\def\oznacenieCelku{Kolekcia učebných textov}

% -----------------------------------------------------------------------------


\def\KUTporadoveCislo{devAS0}

\def\oznacenieVerzie{v1.0}
% \def\oznacenieVerzie{\phantom{v1.0}}

\def\mesiacRok{august 2025}

\def\authorslabel{RJ}






% -----------------------------------------------------------------------------

\begin{document}

% -----------------------------------------------------------------------------
% Uvodny nadpis

\noindent
\parbox[t][18mm][c]{0.3\textwidth}{%
    \raisebox{-0.9\height}{%
        \phantom{.}
\includegraphics[height=18mm]{./COMMONFILES/URKFEIlogo.pdf}%
    }%
}%
\parbox[t][18mm][c]{0.7\textwidth}{%
    \raggedleft

    \sffamily
    \fontsize{16pt}{18pt}
    \fontseries{sbc}
    \selectfont

    \noindent
    \textcolor[rgb]{0.75, 0.75, 0.75}{\textls[25]{\oznacenieCelku}}
}%

\noindent
\parbox[t][16mm][b]{0.5\textwidth}{%
    \raggedright

    \color{Gray}
    \sffamily

    \fontsize{12pt}{12pt}
    \selectfont
    \mesiacRok

    \fontsize{6pt}{10pt}
    \selectfont
    \href{https://github.com/OkoliePracovnehoBodu/KUT}{github.com/OkoliePracovnehoBodu/KUT}

    \fontsize{8pt}{10pt}
    \selectfont
    \authorslabel




}%
\parbox[t][16mm][b]{0.5\textwidth}{%
    \raggedleft

    \sffamily

    \fontsize{6pt}{6pt}
    \selectfont

    \textcolor[rgb]{0.68, 0.68, 0.68}{\oznacenieVerzie}


    \fontsize{14pt}{14pt}
    \selectfont

    \bfseries

    
\includegraphics[height=12pt]{./COMMONFILES/KUT_logo_v0.1.pdf}%
    {%
        \textls[-50]{\KUTporadoveCislo}
    }%
}%

% -----------------------------------------------------------------------------




\vspace{6mm}

% ---------------------------------------------
\sffamily
\bfseries
\fontsize{18pt}{21pt}
\selectfont

\begin{flushleft}
    Laboratórne zariadenie AeroShield:\\ orientačný prehľad
\end{flushleft}

\bigskip

% -----------------------------------------------------------------------------
\normalsize
\normalfont
% -----------------------------------------------------------------------------

\lstset{style=mystyle}










\noindent
\lettrine[lines=1, nindent=1pt, loversize=0.0]{C}{ieľom}
textu je opis laboratórneho zariadenia AeroShield predstavujúceho fyzický model spojitého dynamického systému.


\section{Opis dynamického systému}

AeroShield (daľej len \emph{AS}) je laboratórne zariadenie predstavujúce reálny dynamický systém. Pozostáva z jednoramenného kývadla, malého jednosmerného motora s vrtuľkou na vystupnom hriadeli - uloženého na volnom konci ramena. Následne je výstup snímaní magnetickým rotačným enkóderom AS5600, ktorý zabezpečuje presné meranie uhlovej polohy ramena. Osadený potenciometer, ktorý nerobí nič, iba sa točí a jeho pozícia je snímana pomocou Arduino UNO zariadenia, možno ho využiť na rôzne účely napr. predčasné ukončenie simulácie ak je hodnota signálu > 500.

Systém má jeden vstupný signál a jeden výstupný signál. Výstupný signál je priamo úmerný uhlovej polohe ramena, ktorá je snímaná enkóderom. Vstupný signál nepriamo ovláda napájanie motora cez otvorenie mosfet tranzisotra, ktorý priamo zásobuje napätie na motor.

Polohou potenciometra je v podstate signál, ktorý vie uživateľ použiť na rôzne účeli, no neovplyvňuje systém ako taký, bez toho aby si to uživateľ sám neimplementoval.



\section{Rozsahy a jednotky signálov}

Z opisu predmetného dynamického systému vyplýva, že systém má jeden výstupný signál, jeden vstupný signál a manuálne nastaviteľnú polohu potenciometra.

\begin{itemize}
    \item \emph{Vstupný signál} nadobúda hodnoty v rozmedzí $0$ až $255$.

    \item \emph{Výstupný signál} nadobúda hodnoty v rozsahu \ang{-50} az \ang{210}.

    \item \emph{Signal z potenciometra} nadobúda hodnoty v rozsahu $0$ až $1023$.
\end{itemize}


\begin{center}

    \vspace{-10pt}

    \tabcaption{Rozsahy a jednotky signálov}
    \label{tab:rozsahy_a_jednotky_signalu}

    \lstyle

    \begin{tabular*}{\textwidth}{@{ \extracolsep{\fill}} lll}
        \toprule
        Signál & Rozsah hodnôt & Jednotka \\
        \midrule
        Vstup & $0$ až $255$ & - \\
        Výstup & \ang{-50} až \ang{210} & \si{\degreeCelsius} \\
        Potenciometer & $0$ až $1023$ & - \\
        \bottomrule
    \end{tabular*}


\end{center}







\section{Schematické znázornenie systému}


\begin{center}

    \vbox{%


        \makebox[\textwidth][c]{%
            \includegraphics{AS_system.pdf}
        }

        \figcaption{
            Signály systému AS.
        }
        \label{LMOT_sch1}
    }%vbox

\end{center}


\begin{center}

    \vbox{%


        \makebox[\textwidth][c]{%
            \includegraphics{AS_system_pinout.pdf}
        }

        \figcaption{
            Schéma pripojenia laboratórneho zariadenia AS ku GPIO pinom Arduino UNO R$3$, pričom A$0$ je analógový vstup potenciometra, $\sim$D$5$ je výstup akčného zásahu na logickej úrovni \emph{PWM} (\href{https://docs.arduino.cc/learn/microcontrollers/analog-output/}{Pulse-Width Modulation}), I$2$C slúži na odčítanie uhlovej polohy ramena.
        }
        \label{LMOT_sch2}
    }%vbox

\end{center}



\section{Fotografie}

Zoznam fotografií:

\begin{itemize}
    \item Obr. \ref{IMG_6874_toPDF}: Celkový pohľad na laboratórne zariadenie LMOT.
    \item Obr. \ref{IMG_6922_toPDF}: Motor a tachodynamo.
    \item Obr. \ref{IMG_6941_toPDF}: Elektronické obvody zariadenia.
    \item Obr. \ref{IMG_6931_toPDF}: Predný panel so svorkami a potenciometrom.
\end{itemize}





\noindent
\vbox{%

    \makebox[\textwidth][c]{%
        \includegraphics{IMG_6874_toPDF.jpg}
    }

    \figcaption{
        Celkový pohľad na laboratórne zariadenie LMOT.
    }

    \vspace{12pt}

    \label{IMG_6874_toPDF}
}%vbox



\noindent
\vbox{%

    \makebox[\textwidth][c]{%
        \includegraphics{IMG_6922_toPDF.jpg}
    }

    \figcaption{
        Motor a tachodynamo.
    }

    \vspace{12pt}

    \label{IMG_6922_toPDF}
}%vbox




\noindent
\vbox{%

    \makebox[\textwidth][c]{%
        \includegraphics{IMG_6941_toPDF.jpg}
    }

    \figcaption{
        Elektronické obvody zariadenia.
    }

    \vspace{12pt}

    \label{IMG_6941_toPDF}
}%vbox




\noindent
\vbox{%

    \makebox[\textwidth][c]{%
        \includegraphics{IMG_6931_toPDF.jpg}
    }

    \figcaption{
        Predný panel so svorkami a potenciometrom.
    }

    \vspace{12pt}

    \label{IMG_6931_toPDF}
}%vbox

























































% -----------------------------------------------------------------------------

\end{document}

% -----------------------------------------------------------------------------